
\documentclass[11pt]{article}


\setlength{\oddsidemargin}{.0in}
\setlength{\evensidemargin}{.0in}
\setlength{\textwidth}{6.5in}
\setlength{\topmargin}{-.3in}
\setlength{\headsep}{.20in}
\setlength{\textheight}{9.in}
\raggedbottom


% PACKAGES

\usepackage{url}
\usepackage{hyperref}
\usepackage{amsmath}
\usepackage{amssymb}
\usepackage{graphicx}
\usepackage{float}
\usepackage{caption}
\usepackage{amsthm}
\usepackage{color}
\usepackage{tabu}
\usepackage{cleveref}
\definecolor{green}{rgb}{0,0.5977,0}
\usepackage[linesnumbered,algoruled,boxed,lined]{algorithm2e}

% GENERAL SHORTCUTS

% Common symbols
\newcommand{\rr}{\mathbb R}
\newcommand{\cc}{\mathbb C}
\newcommand{\zz}{\mathbb Z}
\newcommand{\nn}{\mathbb N}
\newcommand{\qq}{\mathbb Q}
\newcommand{\ee}{\mathbb E}
\newcommand{\ff}{\mathbb F}
\newcommand{\zzn}[1]{\zz/{#1}\zz}
\newcommand{\id}{\text{id}}

% Common operations
\newcommand{\inv}{^{-1}}
\newcommand{\abs}[1]{\left|{#1}\right|}
\newcommand{\man}[1]{\left|\left|{#1}\right|\right|_1}
\newcommand{\suchthat}{\ | \ }
\newcommand{\tand}{\txt{ and }}
\newcommand{\normalize}[1]{\frac{#1}{\abs{#1}}}
\newcommand{\normallize}[1]{\frac{#1}{\abs{{\abs#1}}}}
\newcommand{\tth}{^\text{th}}
\newcommand{\genseq}[3]{{#1}_1 {#3} {#1}_2 {#3} \dots {#3} {#1}_{#2}}
\newcommand{\seq}[2]{\genseq{#1}{#2}{,}}
\DeclareMathOperator*{\argmin}{argmin}
\DeclareMathOperator*{\argmax}{argmax}
\newcommand{\pp}[1]{\mathbb{P}\left\{#1\right\}}

% General tools
\newcommand{\twocases}[4]{\begin{cases} #2 & #1 \\ #4 & #3 \end{cases}}
\newcommand{\threecases}[6]{\begin{cases} #2 & #1 \\ #4 & #3 \\ #6 & #5 \end{cases}}
\newcommand{\triad}[5]{\threecases{#1 < #2}{#3}{#1 = #2}{#4}{#1 > #2}{#5}}
\newcommand{\emptybox}[2][\textwidth]{\begingroup\setlength{\fboxsep}{-\fboxrule}\noindent\mbox{\rule{0pt}{#2}}\endgroup}
\newcommand{\leavespace}{{\centering\emptybox[15cm]{10cm}}}
\newcommand{\txt}[1]{\text{#1}}
\newcommand{\fcite}[1]{\texttt{#1.pdf}}
\newcommand{\red}[1]{{\color{red} #1}}

% Align mode tools
\newcommand{\stext}[1]{\ \ \ \ \ \text{(#1)}}
\newcommand{\stextn}[1]{\\&\ \ \ \ \ \ \stext{#1}}
\newcommand{\bcause}[1]{\stext{because ${#1}$}}
\newcommand{\snc}[1]{\stext{since ${#1}$}}
\newcommand{\indhyp}{\stext{by the inductive hypothesis}}
\newcommand{\triineq}{\stext{by the triangle inequality}}
\newcommand{\ub}{\stext{by the union bound}}
\newcommand{\bydef}{\stext{by definition}}
\newcommand{\push}{\\ & \ \ \ \ \ \ \ \ \ \ }
\newcommand{\topindent}{&\hspace{1.335em}}

% Graphics
% \newcommand{\ig}[2]{\begin{figure}[H]\begin{center}\includegraphics[scale = {#1}]{#2.png}\end{center}\end{figure}}
% \newcommand{\igc}[3]{\begin{figure}[H]\begin{center}\includegraphics[scale = {#1}]{#2.png}\caption*{#3}\end{center}\end{figure}}
% \newcommand{\ignc}[3]{\begin{figure}[H]\begin{center}\includegraphics[scale = {#1}]{#2.png}\caption{#3}\end{center}\end{figure}}
% \newcommand{\igm}[2]{\begin{figure}[H]\begin{center}\includegraphics[scale = {#1}]{../Media/#2.png}\end{center}\end{figure}}
% \newcommand{\igcm}[3]{\begin{figure}[H]\begin{center}\includegraphics[scale = {#1}]{../Media/#2.png}\caption*{#3}\end{center}\end{figure}}
% \newcommand{\igncm}[3]{\begin{figure}[H]\begin{center}\includegraphics[scale = {#1}]{../Media/#2.png}\caption{#3}\end{center}\end{figure}}
% \newcommand{\ignsm}[4]{\begin{figure}[H]\begin{center}\includegraphics[scale = {#1}]{../Media/#2.png}\caption[#3]{#4}\end{center}\end{figure}}
% \newcommand{\ip}[2]{\begin{figure}[H]\begin{center}\includegraphics[scale = {#1}]{#2.pdf}\end{center}\end{figure}}
% \newcommand{\ipc}[3]{\begin{figure}[H]\begin{center}\includegraphics[scale = {#1}]{#2.pdf}\caption*{#3}\end{center}\end{figure}}
\newcommand{\ipnc}[4]{\begin{figure}[H]\begin{center}\includegraphics[scale = {#1}]{#2.pdf}\caption{#3}\label{#4}\end{center}\end{figure}}
% \newcommand{\ipm}[2]{\begin{figure}[H]\begin{center}\includegraphics[scale = {#1}]{../Media/#2.pdf}\end{center}\end{figure}}
% \newcommand{\ipcm}[3]{\begin{figure}[H]\begin{center}\includegraphics[scale = {#1}]{../Media/#2.pdf}\caption*{#3}\end{center}\end{figure}}
% \newcommand{\ipncm}[3]{\begin{figure}[H]\begin{center}\includegraphics[scale = {#1}]{../Media/#2.pdf}\caption{#3}\end{center}\end{figure}}
% \newcommand{\ipnsm}[4]{\begin{figure}[H]\begin{center}\includegraphics[scale = {#1}]{../Media/#2.pdf}\caption[#3]{#4}\end{center}\end{figure}}

% Algorithms
\makeatletter 
\g@addto@macro{\@algocf@init}{\SetKwInOut{Parameter}{Parameters}} 
\makeatother

% SUBJECT SPECIFIC SHORTCUTS

% Algebra
\newcommand{\pres}[2]{\langle {#1} \ | \ {#2} \rangle}

% Analysis
\newcommand{\infsum}{\sum_{n = 0}^\infty}
\newcommand{\finsum}{\sum_{n = 0}^N}
\newcommand{\limh}{\lim_{h \to 0}}

% Complex analysis
\newcommand{\real}{\text{Re}}
\newcommand{\imag}{\text{Im}}
\newcommand{\nz}{|z|^2}
\newcommand{\Arg}{\text{Arg}}
\newcommand{\Log}{\text{Log}}
\newcommand{\ddz}{\frac{d}{dz}}

% Complexity theory
\newcommand{\PP}{\mathsf{P}}
\newcommand{\NP}{\mathsf{NP}}
\newcommand{\NPh}{$\NP$-hard }
\newcommand{\NPc}{$\NP$-complete }
\newcommand{\BQP}{\mathsf{BQP}}

% Optimization
\newcommand{\succinctlp}[3]{\begin{tabular}{l l}\textbf{#1} & \begin{tabular}{l}$#2$\end{tabular}\\\textbf{subject to} & \begin{tabular}{l}$#3$\end{tabular}\end{tabular}}
\newcommand{\spmax}[2]{\succinctlp{maximize}{#1}{#2}}
\newcommand{\spmin}[2]{\succinctlp{minimize}{#1}{#2}}
\newcommand{\lp}[3]{\begin{tabular}{l l}\textbf{#1} & \begin{tabular}{l}$#2$\end{tabular}\\\textbf{subject to} & \begin{tabular}{l l}#3\end{tabular}\end{tabular}}
\newcommand{\lpmax}[2]{\lp{maximize}{#1}{#2}}
\newcommand{\lpmin}[2]{\lp{minimize}{#1}{#2}}

% Partition theory
\newcommand{\pofn}[1]{p(n \ | \ \text{#1})}

% p-adic numbers
\newcommand{\pabs}[1]{\left|{#1}\right|_p}
\newcommand{\qp}{\qq_p}
\newcommand{\zp}{\zz_p}

% Quantum mechanics
\newcommand{\bra}[1]{\langle {#1} |}
\newcommand{\ket}[1]{| {#1} \rangle}
\newcommand{\braket}[2]{\langle {#1} | {#2} \rangle}

% Topology
\newcommand{\rp}[1]{\rr\text{P}^{#1}}
\newcommand{\cp}[1]{\cc\text{P}^{#1}}
\newcommand{\Ht}{\widetilde{H}}

% Linear algebra
\newcommand{\mTwoTwo}[4]{\begin{bmatrix}{#1}&{#2}\\{#3}&{#4}\end{bmatrix}}
\newcommand{\mTwoThree}[6]{\begin{bmatrix}{#1}&{#2}&{#3}\\{#4}&{#5}&{#6}\end{bmatrix}}
\newcommand{\mTwoFour}[8]{\begin{bmatrix}{#1}&{#2}&{#3}&{#4}\\{#5}&{#6}&{#7}&{#8}\end{bmatrix}}
\newcommand{\mThreeTwo}[6]{\begin{bmatrix}{#1}&{#2}\\{#3}&{#4}\\{#5}&{#6}\end{bmatrix}}
\newcommand{\mThreeThree}[9]{\begin{bmatrix}{#1}&{#2}&{#3}\\{#4}&{#5}&{#6}\\{#7}&{#8}&{#9}\end{bmatrix}}
\newcommand{\mFourTwo}[8]{\begin{bmatrix}{#1}&{#2}\\{#3}&{#4}\\{#5}&{#6}\\{#7}&{#8}\end{bmatrix}}
\newcommand{\vectTwo}[2]{\begin{bmatrix}{#1}\\{#2}\end{bmatrix}}
\newcommand{\rowTwo}[2]{\begin{bmatrix}{#1}&{#2}\end{bmatrix}}
\newcommand{\vectThree}[3]{\begin{bmatrix}{#1}\\{#2}\\{#3}\end{bmatrix}}
\newcommand{\rowThree}[3]{\begin{bmatrix}{#1}&{#2}&{#3}\end{bmatrix}}
\newcommand{\vectFour}[4]{\begin{bmatrix}{#1}\\{#2}\\{#3}\\{#4}\end{bmatrix}}
\newcommand{\rowFour}[4]{\begin{bmatrix}{#1}&{#2}&{#3}&{#4}\end{bmatrix}}

% Category theory
\newcommand{\inl}{\txt{inl}}
\newcommand{\inr}{\txt{inr}}
\newcommand{\obj}{\txt{obj }}

\theoremstyle{plain}
\newtheorem{theorem}{Theorem}
\newtheorem{lemma}[theorem]{Lemma} 
\newtheorem{proposition}[theorem]{Proposition}
\newtheorem{claim}[theorem]{Claim}
\newtheorem{corollary}[theorem]{Corollary}
\newtheorem{conjecture}[theorem]{Conjecture}
\newtheorem{question}[theorem]{Question}
\newtheorem{globalTheorem}{Theorem}
\theoremstyle{definition}
\newtheorem{definition}[theorem]{Definition}
\newtheorem{example}[theorem]{Example}
\newtheorem{exercise}[theorem]{Exercise}

\numberwithin{theorem}{section}

\usepackage{authblk}

